\section{INTRODUÇÃO}
O projeto Night Heist combina um documento de design elaborado por estudantes com um jogo implementado por agente ChatGPT, segundo a declaração do responsável pelo projeto; o material disponível inclui especificação, código e documentação de desenvolvimento, permitindo examinar a relação entre esses artefatos. Essa configuração desloca o interesse da simples existência de um jogo funcional para a identificação de quais decisões estavam documentadas e de quais detalhes foram estabelecidos durante a implementação.

A literatura recente situa o uso de IA generativa em tarefas de engenharia de software, incluindo priorização de requisitos em contexto educacional, e também problematiza a relação entre assistência automatizada e competência de programadores iniciantes (Queiroz; Lima, 2025; Prather et al., 2024). Esses antecedentes justificam examinar artefatos e processos de especificação sem assumir que a produção final, isoladamente, demonstre aprendizagem ou domínio da programação (Prather et al., 2024).

O problema concreto deste caso consiste na distância entre enunciados do GDD e parâmetros executáveis: o documento descreve retorno ao ponto inicial, temporizador e terminal de hack, enquanto o código precisa definir condições espaciais, valores temporais e efeitos dessas interações. Sem rastreabilidade, uma decisão adicionada para tornar a regra executável pode ser confundida com requisito original, e uma divergência pode ser interpretada como defeito sem consulta ao contexto de design.

A pergunta de pesquisa é: \emph{como as decisões explicitadas no GDD discente se relacionam com a implementação do Night Heist e quais correspondências, divergências e lacunas de especificação podem ser documentadas por inspeção estática?} A formulação delimita o objeto à relação entre artefatos, conforme o protocolo do projeto, e não à avaliação individual dos estudantes.

O objetivo geral é analisar essa relação mediante uma matriz de rastreabilidade com critérios observáveis, localização das evidências e distinção entre estados de verificação. Os objetivos específicos são extrair unidades do GDD, relacioná-las ao código, classificar sua correspondência, identificar parâmetros não especificados e registrar os limites de atribuição das decisões. Neste manuscrito, esses objetivos são operacionalizados por um piloto de seis unidades, sem apresentar o piloto como inventário completo do documento.

A contribuição apresentada é um procedimento documental explicitado e acompanhado de exemplos auditáveis, que pode orientar a revisão humana do projeto e a continuidade da análise. Sua relevância científica reside em tornar examinável a tradução de especificação humana para implementação assistida por IA; a originalidade em relação a outros procedimentos e a utilidade em outros projetos ainda precisam ser investigadas, sem alegação de método validado ou superioridade da IA (Queiroz; Lima, 2025).

\section{FUNDAMENTAÇÃO E DELIMITAÇÃO CONCEITUAL}
\subsection{GDD como especificação do objeto de análise}
No documento fornecido, o GDD organiza visão geral, estrutura do cenário, fluxo de missão, mecânicas, arquitetura, recursos visuais e backlog; seus enunciados têm níveis distintos de detalhe. O cenário de quatro pavimentos e o limiar de exposição à câmera oferecem critérios diretos de comparação, enquanto a referência ao terminal de hack não estabelece, no texto lido, sua duração nem todos os seus efeitos.

Essa heterogeneidade exige separar a presença de uma ideia de jogo da especificação suficiente de sua regra operacional: citar um temporizador não determina sua duração, e mencionar um retorno ao ponto inicial não resolve automaticamente a tolerância espacial. No estudo, a unidade de análise é um enunciado de decisão ou requisito associado a um critério verificável; um mesmo requisito amplo pode exigir mais de uma unidade para distinguir componentes e parâmetros.

O antecedente de uso do ChatGPT na priorização de requisitos oferece um vínculo temático com decisões de engenharia de software, mas seus resultados não são transferidos para este jogo nem utilizados como validação da matriz proposta (Queiroz; Lima, 2025). A matriz é um instrumento operacional do caso, cuja clareza e consistência ainda dependem de aplicação e revisão humana documentadas.

\subsection{Correspondência, ambiguidade e proveniência}
Correspondência designa a relação entre critério explícito e evidência examinada; ela não funciona, neste estudo, como índice global de qualidade, diversão ou valor pedagógico. Uma regra pode divergir literalmente da especificação e ainda ser adequada para a interação, hipótese que só pode ser examinada com uma justificativa documentada ou verificação específica.

Ambiguidade e insuficiência de especificação são registradas separadamente da categoria de correspondência, porque a ausência de parâmetro numérico não equivale à comprovação de uma implementação incorreta. A categoria não verificável conserva essa insuficiência como resultado documental, em vez de converter silenciosamente a escolha do código em requisito do GDD.

Proveniência da decisão é a identificação do registro que sustenta quem resolveu uma ambiguidade ou definiu um parâmetro; o relato geral de programação por agente não identifica o autor de cada escolha. Na ausência de conversas originais, solicitações ou decisões assinadas, a matriz conserva o responsável como desconhecido, inclusive nos casos em que o código contém um valor numérico claro.

\subsection{Artefato técnico e evidência educacional}
Prather et al. (2024) examinam benefícios e danos da IA generativa para programadores iniciantes, oferecendo fundamento para cautela ao associar desempenho em tarefas assistidas a competência individual. Por isso, a existência do jogo e do GDD é tratada aqui como evidência documental da experiência relatada, sem medida de aprendizagem, comparação entre estudantes ou conclusão sobre eficácia pedagógica (Prather et al., 2024).

Lelli, Santos e Castelo Branco (2024) apresentam práticas gamificadas no ensino de verificação e validação, fornecendo um antecedente educacional próximo ao interesse pelo desenvolvimento de jogos. O vínculo temático não torna o Night Heist uma intervenção educacional avaliada: o corpus atual sustenta a análise técnica de decisões e não resultados de ensino (Lelli; Santos; Castelo Branco, 2024).

\section{MÉTODO}
\subsection{Desenho e escopo}
Foi realizada análise documental piloto de caso único com inspeção estática de código, conforme o protocolo documental descrito nesta seção. A escolha do caso decorre da disponibilidade do GDD e do repositório; as seis unidades foram selecionadas intencionalmente para exercitar correspondência explícita, divergência literal e falta de parâmetro, sem procedimento probabilístico ou pretensão de representatividade do documento.

O piloto permite verificar se as categorias produzem registros compreensíveis e localizáveis, mas não estima a frequência de problemas no jogo inteiro. Não houve questionário, entrevista, observação sistemática de estudantes, comparação de aprendizagem ou coleta de notas no corpus descrito. A eventual ampliação para essas dimensões configurará outro escopo e precisará de protocolo próprio, em lugar de inferência retrospectiva a partir deste material (Brasil, 2016).

\subsection{Corpus e controle de versões}
O GDD e o código são identificados abaixo como fontes primárias do corpus. Os registros U01--U06 são os localizadores dos resultados, sem autocitação bibliográfica.

O corpus primário reúne a leitura textual do PDF fornecido pelo usuário e os arquivos do baseline identificado pelo commit \nolinkurl{72699f0c2752d029d8e8c3b48420c18a70dbe2a1}. A identificação por commit fixa o objeto técnico e evita que mudanças posteriores sejam incorporadas silenciosamente à análise. Os instrumentos científicos, por sua vez, são documentos de trabalho da branch \texttt{docs/projeto-academico-cientifico}, proposta na PR 2, e não arquivos integrantes desse baseline do jogo.

O arquivo no Drive é denominado \emph{GDD -- Night Heist Escape Phase (Versão 0).pdf}, enquanto o conteúdo identifica \emph{Night Heist: Escape Phase V2}, versão 0.2, MVP sem grade descendente; as duas identificações foram preservadas no registro de fonte. Os metadados consultados informaram tamanho de 142006 bytes e modificação em 8 de outubro de 2026, mas os bytes originais não foram materializados no ambiente de análise.

Consequentemente, não foi calculado hash do PDF nem validada sua paginação, e a inspeção visual do mockup permanece pendente. Os localizadores utilizados são seções e identificadores de mecânica, adequados à leitura textual disponível, sem inventar números de página. O artigo incorpora o GDD como fonte e objeto de análise, não como reprodução integral anexada.

A documentação interna de design, backlog e validação foi tratada como fonte auxiliar do repositório, sem substituir o GDD original. Conversas com o agente e registros que individualizem decisões não compõem o corpus atual, motivo pelo qual autoria de parâmetros não foi deduzida a partir do código ou do relato geral da experiência.

\subsection{Instrumento e categorias}
A matriz registra identificador da unidade, origem, localizador e trecho do GDD, critério, localizador da implementação, teste associado quando existente, evidência, categoria, ambiguidade, responsável, fonte da decisão, estado técnico e revisão humana. A presença de coluna de teste não significa execução: nos seis registros iniciais, o estado é inspeção estática e a revisão humana está pendente.

\begin{singlespace}
\begin{longtable}{p{0.22\textwidth}p{0.69\textwidth}}
\caption{Categorias operacionais de correspondência.}\label{tab:categories}\\
\toprule Categoria & Regra adotada no instrumento \\\midrule\endfirsthead
\toprule Categoria & Regra adotada no instrumento \\\midrule\endhead
Preservado & Todos os critérios explícitos da unidade são atendidos na evidência examinada.\\
Parcial & Parte dos critérios é atendida; as partes não atendidas precisam ser identificadas.\\
Alterado & Há implementação correspondente que contradiz pelo menos um critério explícito.\\
Omitido & Unidade explícita aplicável não tem implementação após busca e verificação registradas.\\
Adicionado & Há comportamento implementado sem unidade correspondente no GDD completo.\\
Não verificável & Fonte, escopo ou evidência são insuficientes para decidir a correspondência.\\
\bottomrule
\multicolumn{2}{l}{\footnotesize Fonte: codebook operacional do projeto .}\\
\end{longtable}
\end{singlespace}

O instrumento distingue as ambiguidades lexical, conflitante, falta de parâmetro, escopo incerto e desconhecida, além de ausência de ambiguidade. Para atribuição de responsável, admite discente, orientador, agente, conjunto ou desconhecido, sempre condicionado a registro original verificável. As categorias são propostas para este caso e não constituem escala psicométrica ou instrumento previamente validado.

\subsection{Procedimento de inspeção}
Primeiro, foram identificadas unidades explícitas e lacunas numéricas na leitura textual do GDD; depois, cada unidade foi associada a um critério observável antes da localização correspondente no código. A comparação conservou a diferença entre estrutura declarada e comportamento condicionado, de modo que uma constante foi usada como evidência estrutural e uma condição como evidência de regra lógica.

Em seguida, foram examinados \texttt{scripts/level.gd}, \texttt{scripts/mission.gd} e \texttt{scripts/main.gd}, registrando símbolo ou método associado a cada unidade. A classificação foi acompanhada de justificativa, sem deduzir que um trecho encontrado foi necessariamente executado em todas as plataformas previstas. A última etapa do piloto consistiu em consolidar a matriz e seus limites; revisão independente e testes em runtime permanecem etapas futuras.

\subsection{Ética, transparência e condições editoriais}
A unidade de análise é a decisão documentada relacionada ao código; a identificação de autoria foi fornecida pelo usuário, sem atribuição individual de desempenho ou de cada decisão. O uso de produções escolares exige avaliar direitos e enquadramento pertinente, e este manuscrito não declara dispensa automática de avaliação ética apenas pela ausência de questionários (Brasil, 2016).

A concepção do recorte, a organização do instrumento, a inspeção documental e a redação deste manuscrito receberam assistência de ChatGPT/Codex; os resultados permanecem sujeitos à revisão intelectual dos pesquisadores humanos. A autorização de continuidade do usuário suspendeu, nesta etapa de elaboração, a adaptação às condições da revista anteriormente considerada; não há declaração de conformidade editorial nem submissão realizada. O repositório não possui licença geral do código no baseline auditado, portanto acesso público não foi apresentado como comprovação de licença open source do conjunto.

\section{RESULTADOS DOCUMENTAIS PRELIMINARES}
\subsection{Correspondências explícitas}
A unidade U01 relaciona a descrição de quatro pavimentos, na seção 2 do GDD, à lista \texttt{FLOORS} de \texttt{scripts/level.gd}; a lista contém quatro valores, sustentando a classificação preservado para a cardinalidade da estrutura. Esse resultado não avalia disposição visual, acessibilidade entre pavimentos ou fidelidade de todos os elementos do cenário, pois esses aspectos não integram o critério da unidade.

Na unidade U02, a regra M-01 estabelece falha por esgotamento do temporizador, e o método \texttt{tick} de \texttt{scripts/mission.gd} termina a missão sem sucesso quando o tempo restante chega a zero. A correspondência estática sustenta a classificação preservado para a condição lógica inspecionada, sem confirmar por execução a atualização do temporizador ou sua apresentação na interface.

Na unidade U03, a regra M-03 fixa exposição superior a 0,5 segundo como gatilho de alerta, e \texttt{sensor\_tick} contém a condição de exposição superior a esse limiar. A coincidência do valor e do operador sustenta a classificação preservado, mas a leitura estática não verifica geometria do cone, detecção de colisões ou estabilidade temporal em execução.

\subsection{Divergência literal no retorno}
A unidade U04 compara a referência ao retorno ao ponto exato de início, nas seções 1.1 e 3.1, com a condição espacial utilizada em \texttt{main.gd} e encaminhada à missão. O código aceita distância inferior a 30 pixels, enquanto o enunciado literal do GDD não explicita esse raio; o registro foi classificado como alterado.

A classificação identifica uma diferença entre especificação textual e regra implementada, sem concluir que a tolerância seja inadequada. A expressão ponto exato pode ter sido empregada de forma informal pelo designer, e um raio pode constituir uma interpretação aceitável, mas nenhuma dessas hipóteses foi tratada como decisão confirmada no corpus. Uma revisão humana poderá manter a classificação literal, reconhecer ambiguidade lexical ou registrar mudança aprovada, desde que preserve a evidência e a justificativa da revisão.

\subsection{Parâmetros sem especificação suficiente}
A unidade U05 examina a duração da missão: o GDD prevê temporizador, mas não fornece duração numérica no trecho lido; o código adota 180 segundos. A fidelidade numérica foi classificada como não verificável por falta de parâmetro, pois não existe valor explícito de referência para confirmar ou contradizer o valor implementado.

A unidade U06 examina o terminal de hack mencionado no pavimento 2: o código desativa a segurança por 12 segundos, mas o GDD lido não parametriza esse efeito e sua duração. Também nesse caso, o registro não verificável expressa insuficiência de especificação, e não ausência de implementação ou comprovação de erro.

A distinção entre uma função mencionada e seus parâmetros é central nesses dois registros: o comportamento não é inteiramente alheio ao GDD, mas contém concretizações que o texto disponível não permite comparar numericamente. Por isso, o piloto não classificou automaticamente essas unidades como adicionadas, categoria que exigiria verificar a ausência de correspondente no GDD completo.

\begin{singlespace}
\begin{longtable}{p{0.07\textwidth}p{0.25\textwidth}p{0.34\textwidth}p{0.20\textwidth}}
\caption{Síntese das seis unidades selecionadas.}\label{tab:results}\\
\toprule ID & GDD e critério & Evidência estática & Classificação \\\midrule\endfirsthead
\toprule ID & GDD e critério & Evidência estática & Classificação \\\midrule\endhead
U01 & Seção 2: quatro pavimentos & \texttt{level.gd:FLOORS}: quatro valores & Preservado\\
U02 & M-01: tempo zerado implica falha & \texttt{mission.gd:tick}: termina sem sucesso & Preservado\\
U03 & M-03: exposição superior a 0,5 s & \texttt{level.gd:sensor\_tick}: limiar correspondente & Preservado\\
U04 & Seções 1.1 e 3.1: ponto exato inicial & \texttt{main.gd}: distância inferior a 30 px & Alterado, leitura literal\\
U05 & M-01: temporizador sem duração numérica & \texttt{mission.gd:DURATION}: 180 s & Não verificável\\
U06 & Seção 2: terminal de hack sem efeito parametrizado & \texttt{level.gd:\_hack}: segurança desativada por 12 s & Não verificável\\
\bottomrule
\multicolumn{4}{p{0.93\textwidth}}{\footnotesize Fonte: GDD  e matriz do projeto . Todos os registros: inspeção estática, responsável desconhecido e revisão humana pendente.}\\
\end{longtable}
\end{singlespace}

Não foram calculados percentuais globais de fidelidade ou qualidade, pois as seis unidades constituem seleção intencional e não o conjunto de requisitos do GDD. O quadro apresenta ocorrências do piloto e não permite concluir ausência das demais categorias no jogo completo.

\section{DISCUSSÃO}
\subsection{O que a rastreabilidade permite observar}
Os registros U01--U03 mostram que critérios relativamente delimitados podem receber correspondência documental clara, desde que o resultado seja restrito ao aspecto efetivamente inspecionado. O achado útil é a possibilidade de apontar onde a cardinalidade, a condição de falha e o limiar aparecem no código, evitando uma avaliação genérica de que o jogo seguiu o GDD.

O registro U04 demonstra que uma condição aparentemente simples exige explicitar a interpretação adotada pelo analista: igualdade espacial e proximidade são regras diferentes, embora possam cumprir a mesma intenção de design. Preservar o trecho do GDD, o critério literal e o limiar do código permite revisar a classificação sem apagar a evidência inicial.

Os registros U05 e U06 tornam visível o limite do documento para verificar parâmetros implementados; a categoria não verificável impede afirmar fidelidade numérica quando o número não foi especificado. Esse resultado orienta a gestão de requisitos do próprio projeto, porque identifica perguntas concretas a serem resolvidas e registradas, como duração pretendida e alcance da desativação da segurança.

\subsection{Implicações para design e gestão do desenvolvimento}
No contexto deste caso, o GDD pode ser revisto para distinguir regras essenciais, parâmetros negociáveis e exemplos de implementação, preservando a versão original como fonte histórica. Essa recomendação decorre dos casos observados, sobretudo da diferença entre regra funcional e parametrização, sem alegar eficácia geral de um formato de GDD.

A gestão do desenvolvimento pode utilizar uma decisão documentada por alteração, contendo requisito de origem, motivo, responsável comprovado e versão afetada; o instrumento já contém campos compatíveis com esse registro. Um raio de retorno aprovado, por exemplo, passaria a ter uma justificativa explícita em vez de permanecer apenas como valor no código.

A literatura sobre priorização de requisitos com ChatGPT oferece um antecedente para examinar a participação da IA em atividades de engenharia de software, mas o presente piloto focaliza correspondência pós-implementação e não priorização (Queiroz; Lima, 2025). Essa delimitação permite formular uma contribuição estreita e verificável, sem reivindicar uma nova metodologia educacional a partir de um único caso.

\subsection{Relevância educacional possível e seus limites}
A matriz pode integrar uma atividade de leitura de requisitos, comparação com código e revisão de decisões, porque organiza questões concretas do projeto, mas essa possibilidade constitui proposta de uso e não resultado de aprendizagem. A análise poderia apoiar discussão sobre responsabilidade de design e justificativa de parâmetros sem exigir, no escopo atual, medir a competência individual dos autores.

O cuidado em separar proposta de resultado é compatível com as questões levantadas por Prather et al. (2024) sobre assistência generativa e programadores iniciantes. Para sustentar uma conclusão educacional futura, seriam necessários desenho, evidências e enquadramento específicos; o jogo pronto e a matriz documental não substituem essas condições (Prather et al., 2024; Brasil, 2016).

\section{LIMITAÇÕES E CONTINUIDADE VERIFICÁVEL}
A primeira limitação é o corpus parcial em termos materiais: foi obtida leitura textual do GDD, mas não sua cópia binária para hash, conferência de páginas e inspeção visual. A versão interna divergente do nome do arquivo também exige confirmação humana, sem presumir equivalência entre todas as versões mencionadas.

A segunda limitação é a inspeção estática: os resultados indicam estruturas e condições no código, mas não demonstram sua manifestação em runtime. Os arquivos de teste existentes no repositório são artefatos consultáveis, e sua presença não comprova execução recente, êxito em plataformas específicas ou cobertura das seis unidades.

A terceira limitação é analítica: seis unidades selecionadas não representam todo o GDD, e as categorias ainda não foram revisadas por pesquisadores humanos. O instrumento não possui avaliação de concordância e não sustenta estimativas de confiabilidade entre avaliadores; qualquer medida futura deverá informar procedimento, unidades, denominador e divergências observadas.

A quarta limitação é de proveniência: faltam registros originais que identifiquem quem decidiu cada parâmetro, o que impede separar contribuições individuais de alunos, orientador e agente. A assistência de IA à própria análise também deve ser considerada na revisão humana, especialmente na escolha de unidades, interpretação de ambiguidades e redação das conclusões.

A continuidade proposta consiste em conferir os bytes e a identidade da fonte, revisar humanamente as seis classificações, ampliar a extração às demais regras do GDD e executar verificações técnicas com versão do ambiente e resultados registrados. Essas etapas deverão conservar a distinção entre evidência documental, execução técnica e eventual avaliação educacional, para que mudanças de escopo não sejam apresentadas como resultados do piloto atual.

\section{CONSIDERAÇÕES FINAIS}
O piloto documentou correspondências estáticas entre GDD e código em três unidades selecionadas, uma divergência literal de condição espacial e dois casos em que a especificação não permite comparação numérica. Esses resultados respondem parcialmente à pergunta de pesquisa, restritos aos critérios e localizadores apresentados, sem evidência de cobertura integral ou execução recente.

A contribuição concreta é a matriz rastreável acompanhada de categorias, justificativas e estados de revisão, que torna explícita a diferença entre fidelidade verificável, interpretação de design e insuficiência documental. O procedimento constitui uma proposta aplicada ao caso, cuja consistência, replicabilidade em outros projetos e utilidade precisam de revisão e investigação adicionais.

A experiência relatada pode fundamentar um trabalho científico documental sobre design e gestão do desenvolvimento assistido por IA, desde que o artigo permaneça ancorado nos artefatos e em suas limitações. Os achados não demonstram aprendizagem dos discentes, qualidade global do jogo ou superioridade da implementação por IA; a conclusão científica alcançada é a possibilidade de examinar e registrar a tradução de decisões com uma cadeia explícita de evidências (Prather et al., 2024).
